% Part: sets-functions-relations
% Chapter: relations
% Section: special-properties

\documentclass[../../../include/open-logic-section]{subfiles}

\begin{document}

\olfileid{sfr}{rel}{prp}
\olsection{تعلقاں دیاں خاص خاصیتاں (Special Properties of Relations)}

\begin{intro}
تعلقاں دیاں کجھ قسماں اینیاں عام نیں کہ اوہناں نوں خاص ناں دتے گئے نیں۔
مثال دے طور تے، $\le$ تے $\subseteq$ دونویں اپنے اپنے دائرے دے عنصراں نوں
رلدے ملدے طریقیاں نال جوڑدے نیں (مثلاً $\Nat$، جدوں تعلق $\le$ ہووے،
تے $\Pow{A}$، جدوں تعلق $\subseteq$ ہووے)۔ ایہ ٹھیک ٹھیک سمجھن لئی کہ
ایہ تعلق کِویں رلدے ملدے نیں تے کِویں وکھرے نیں، اسی اوہناں دی ونڈ
کجھ خاص خاصیتاں دے حساب نال کردے آں جیہڑیاں تعلقاں وچ ہو سکدیاں نیں۔
ایہناں وچوں کجھ خاصیتاں (دے میل) خاص اہم نکلدے نیں: ترتیباں (orders) تے
ہم ارزیت دے تعلق (equivalence relations)۔
\end{intro}

\begin{defn}[انعکاسیت (Reflexivity)]
اک تعلق $R \subseteq A^2$ اُدوں تے صرف اُدوں \emph{انعکاسی (reflexive)} ہوندا اے جدوں
ہر $x \in A$ لئی $Rxx$ ہووے۔
\end{defn}

\begin{defn}[تعدیت (Transitivity)]
اک تعلق $R \subseteq A^2$ اُدوں تے صرف اُدوں \emph{متعدی (transitive)} ہوندا اے جدوں،
جد وی $Rxy$ تے $Ryz$ ہون، تاں $Rxz$ وی ہووے۔
\end{defn}

\begin{defn}[تناظر (Symmetry)]
اک تعلق $R \subseteq A^2$ اُدوں تے صرف اُدوں \emph{متناظر (symmetric)} ہوندا اے جدوں،
جد وی $Rxy$ ہووے، تاں $Ryx$ وی ہووے۔
\end{defn}

\begin{defn}[ضد تناظر (Anti-symmetry)]
اک تعلق $R \subseteq A^2$ اُدوں تے صرف اُدوں \emph{ضد تناظری (anti-symmetric)} ہوندا اے جدوں،
جد وی $Rxy$ تے $Ryx$ دونویں ہون، تاں $x=y$ ہووے (یا ہور لفظاں وچ:
جے $x\neq y$ ہووے، تاں یا $\lnot Rxy$ یا $\lnot Ryx$ ہووے)۔
\end{defn}

\begin{explain}
متناظر تعلق وچ $Rxy$ تے $Ryx$ ہمیشہ اکٹھے درست ہوندے نیں، یا دونویں
درست نہیں ہوندے۔ ضد تناظری تعلق وچ $Rxy$ تے $Ryx$ اکٹھے صرف اوہدوں
درست ہو سکدے نیں جدوں $x = y$ ہووے۔ خیال رکھو کہ ایہ اوہناں دا
اکٹھے درست ہونا \emph{لازم} نہیں کردا: $Rxy$ تے $Ryx$ دا $x = y$ ہون اُتے
درست ہونا صرف خارج نہیں کیتا گیا۔ سو ضد تناظری تعلق انعکاسی ہو سکدا
اے، پر ایہ گل نہیں کہ ہر ضد تناظری تعلق انعکاسی ہووے۔ ایہ وی یاد رکھو
کہ ضد تناظری ہونا تے صرف متناظر نہ ہونا دو وکھریاں شرطاں نیں۔ درحقیقت،
اک تعلق اکو ویلے متناظر تے ضد تناظری دونویں ہو سکدا اے (مثلاً اکو ہون
دا تعلق ایہو جیہا اے)۔
\end{explain}

\begin{defn}[باہمی موازنہ پذیری (Connectivity)]
اک تعلق $R \subseteq A^2$ \emph{باہم قابلِ موازنہ (connected)} ہوندا اے جے ہر $x,y\in
A$ لئی، جے $x \neq y$، تاں یا $Rxy$ یا $Ryx$ ہووے۔
\end{defn}

\begin{prob}
ایہو جہے تعلقاں دیاں مثالاں دیو جیہڑے (الف) انعکاسی تے متناظر ہون پر
متعدی نہ ہون، (ب) انعکاسی تے ضد تناظری ہون، (ج) ضد تناظری تے متعدی
ہون پر انعکاسی نہ ہون، تے (د) انعکاسی، متناظر تے متعدی ہون۔ عدداں یا
سیٹاں اُتے تعلق نہ ورتو۔
\end{prob}
 
\begin{defn}[اپنے آپ نال تعلق نہ ہونا (Irreflexivity)]
اک تعلق $R \subseteq A^2$ \emph{اپنے آپ نال تعلق نہ ہون والا (irreflexive)} آکھیا جاندا
اے جے ہر $x \in A$ لئی $Rxx$ نہ ہووے۔
\end{defn}

\begin{defn}[یک طرفہ ہونا (Asymmetry)]
اک تعلق $R \subseteq A^2$ \emph{یک طرفہ (asymmetric)} آکھیا جاندا اے جے کسے وی جوڑے
$x,y\in A$ لئی $Rxy$ تے $Ryx$ دونویں اکٹھے نہ ہون۔
\end{defn}

خیال رکھو کہ جے $A \neq \emptyset$ ہووے، تاں $A$ اُتے اپنے آپ نال تعلق
نہ ہون والا کوئی وی تعلق انعکاسی نہیں ہوندا، تے $A$ اُتے ہر یک طرفہ تعلق
ضد تناظری وی ہوندا اے۔ پر ایہو جہے $R \subseteq A^2$ موجود نیں جیہڑے
انعکاسی وی نہیں تے ہر عنصر دا اپنے آپ نال بے تعلق ہونا وی لازم نہیں کردے؛ تے ایہو جہے ضد تناظری
تعلق وی نیں جیہڑے یک طرفہ نہیں۔

\end{document}
